\section*{第 \ref{ch:b1:gas-exchange} 章——呼吸性气体交换}

\begin{solution}{exo:b1:gas-exchange:1}
$\dot V = K(A/x)(P_1 - P_2)$：面积更大，屏障更薄，
外侧\omterm{def:b1:gas-exchange:partial}{分压}更高（通气），内侧\omterm{def:b1:gas-exchange:partial}{分压}更低
（灌注）。
\end{solution}

\begin{solution}{exo:b1:gas-exchange:2}
每升含氧量低三十倍（\qty{7}{} 对 \qty{210}{mL/L}）；
密度高八百倍，黏度高五十倍，因此泵送
代价高昂；氧气在其中的扩散速度慢一万倍，因此
介质必须横跨整个表面移动。
\end{solution}

\begin{solution}{exo:b1:gas-exchange:3}
并流：水与血液均以 \qty{12}{kPa} 离开；
提取 $(21 - 12)/21 = \qty{43}{\%}$。\omterm{def:b1:gas-exchange:gill}{逆流}：水
以 \qty{4}{kPa} 离开，血液为 \qty{18}{kPa}；提取 $(21 - 4)/21 =
\qty{81}{\%}$。
\end{solution}

\begin{solution}{exo:b1:gas-exchange:4}
溶解的（\qty{7}{\%}）、与血红蛋白氨基结合的（\qty{23}{\%}）、
以碳酸氢根形式存于\omterm{def:b1:mammal-organization:compartments}{血浆}中（\qty{70}{\%}）。
\end{solution}

\begin{solution}{exo:b1:gas-exchange:5}
每分通气量 \qty{6.75}{L/min}；\omterm{def:b1:gas-exchange:lung}{肺泡}通气量 $15\times 0.3 = \qty{4.5}{L/min}$；
摄氧量 $4.5\times 0.07 = \qty{0.32}{L/min}$。浅呼吸：
每分通气量仍为 \qty{6.75}{L/min}，但\omterm{def:b1:gas-exchange:lung}{肺泡}通气量为 $30\times 0.075 =
\qty{2.25}{L/min}$，摄氧量为 \qty{0.16}{L/min}——仅为一半，
呼吸做功相同，但无效腔每次呼吸都要付出代价。
\end{solution}

\begin{solution}{exo:b1:gas-exchange:6}
$2/(6\times 0.75) = \qty{0.44}{L}$，每小时，相当于其
体积 \qty{20}{mL} 的二十二倍。
\end{solution}

\begin{solution}{exo:b1:gas-exchange:7}
$A/x$ 降低 $2\times 4 = 8$ 倍，为 \qty{31}{mL/min}。要恢复
到 \qty{250}{}，压差须达 \qty{64}{kPa}——不可能，因为
\omterm{def:b1:gas-exchange:lung}{肺泡}压在空气中不可超过 \qty{13}{kPa}，静脉血
的压强不可低于零；患者只能通过
吸氧提高，而且即便如此也达不到八倍。
\end{solution}

\begin{solution}{exo:b1:gas-exchange:8}
\omterm{prop:b1:gas-exchange:transport}{碳酸酐酶}位于\omterm{def:b1:cell-unit-of-life:cell}{细胞}内，因此碳酸氢根在
那里生成并积累；它顺浓度梯度经由一种
交换体扩散出去，该交换体每排出一个碳酸氢根便引入一个氯离子，
以保持电荷平衡。若无此交换，碳酸氢根将连同其质子一并
滞留于\omterm{def:b1:cell-unit-of-life:cell}{细胞}内，抬高\omterm{def:b1:cell-unit-of-life:cell}{细胞}渗透压，水分随之
进入：\omterm{def:b1:cell-unit-of-life:cell}{细胞}会在\omterm{def:b1:body-plans-tissues:tissue}{组织}中肿胀（现状下已略有肿胀
）。
\end{solution}

\begin{solution}{exo:b1:gas-exchange:9}
呼吸冲动源于升高的 $\mathrm{CO_2}$。过度通气
使 $\mathrm{CO_2}$ 远低于正常值，却几乎不增加氧气
（血红蛋白原已饱和）；潜水时氧气降至
失去意识的水平，而 $\mathrm{CO_2}$ 尚未回升到
迫使人吸气的水平。若无过度通气，
$\mathrm{CO_2}$ 达到该水平时氧气仍然充足，于是
潜水者浮出水面。
\end{solution}

\begin{solution}{exo:b1:gas-exchange:10}
$A = \pi(10^{-5})^2 = \qty{3.1e-10}{m^2}$；$\Delta c = JL/DA =
10^{-9}\times 10^{-3}/(2\times 10^{-5}\times 3.1\times 10^{-10}) =
\qty{160}{mol/m^3}$——远大于可用的 \qty{8.7}{}，因此这样细的管
道仅靠扩散无法在 \qty{1}{mm} 处为肌肉供氧；
管径较宽（\qty{100}{\micro m}：$\Delta c = \qty{6.4}{mol/m^3}$，
勉强可行），否则需要主动通气。距离为 \qty{10}{mm} 时
需求又增至十倍。由于管长随体型增大，而需求
随体积增大，因此经气管的扩散
将昆虫体型限制在几厘米，更大只能靠泵送。
\end{solution}

\begin{solution}{exo:b1:gas-exchange:11}
鱼类：水单向流动，血液与之反向（\omterm{def:b1:gas-exchange:gill}{逆流}），
提取 \qty{80}{\%}。鸟类：空气单向穿过
副支气管，血液成直角流动（交叉流），提取
率居中，但高于哺乳动物，交换面
从不接触陈旧空气。潮式通气的肺将新鲜空气与
残气混合，因此交换面至多只能接触\omterm{def:b1:gas-exchange:lung}{肺泡}
混合气，血液也只能与之达到平衡：无论
交换面如何，动脉血压都不可超过\omterm{def:b1:gas-exchange:lung}{肺泡}压，而
\omterm{def:b1:gas-exchange:lung}{肺泡}压便无法接近吸入压，只要残气量仍存在
。
\end{solution}

\begin{solution}{exo:b1:gas-exchange:12}
$\mathrm{CO_2}$ 是更好的信号，因为其动脉水平
对通气的反应直接而陡峭（通气减半时其
水平即翻倍），经 \omterm{def:b1:water-small-molecules:ph}{pH} 由中枢感受器精确测得，
而氧气位于血红蛋白曲线的平坦顶部，
在正常通气范围内变化很小，因此
$\mathrm{CO_2}$ 恒温器附带地使氧保持适当水平。它失效的情形是
当氧气独立于 $\mathrm{CO_2}$ 下降时——例如在高海拔，
吸入氧气低而 $\mathrm{CO_2}$ 产生量并未降低，
于是 $\mathrm{CO_2}$ 信号提示“少呼吸”，而氧气提示
“多呼吸”；只有氧气严重不足时，颈动脉体
才会覆盖此信号，而这种折中（伴碱中毒的过度通气）
需要数天的肾脏调节才能稳定。
\end{solution}

\begin{solution}{pb:b1:gas-exchange:1}
\textbf{1.} 空气 $210/22.4 = \qty{9.4}{mmol/L}$；水 $7/22.4 =
\qty{0.31}{mmol/L}$。
\textbf{2.} 30。
\textbf{3.} 空气：$1.2/9.4 = \qty{0.13}{g}$；水：$1000/0.31 =
\qty{3200}{g}$——质量多 \num{25000} 倍。
\textbf{4.} 需水 $210/7 = \qty{30}{L}$。
\textbf{5.} 水过重且含氧过少，难以往复推动；
水须只流过一次，这也使\omterm{def:b1:gas-exchange:gill}{逆流}成为可能；空气足够廉价，
可以往复两次。
\textbf{6.} 每分通气量 $12\times 0.5 = \qty{6}{L/min}$；\omterm{def:b1:gas-exchange:lung}{肺泡}
通气量 $12\times 0.35 = \qty{4.2}{L/min}$。
\textbf{7.} $150/500 = \qty{30}{\%}$，占每次呼吸量及每分通气量的比例。
\textbf{8.} 送达 $4.2\times 0.21 = \qty{0.88}{L/min}$；
摄氧量 $4.2\times 0.07 = \qty{0.29}{L/min}$，接近 \qty{250}{mL/min}。
\textbf{9.} $0.25/(6\times 0.21) = \qty{20}{\%}$（相当于
到达\omterm{def:b1:gas-exchange:lung}{肺泡}部分的 \qty{28}{\%}）。
\textbf{10.} $0.14\times(101 - 6.3) = \qty{13.3}{kPa}$。
\textbf{11.} $250/5 = \qty{50}{mL/L}$；静脉血 \qty{150}{mL/L}，
\qty{75}{\%} 饱和。
\textbf{12.} $3000/20 = \qty{150}{mL/L}$；静脉血 \qty{50}{mL/L}，
\qty{25}{\%} 饱和。
\textbf{13.} $50/60 = \qty{0.83}{mL/min}$。
\textbf{14.} $0.83/(7\times 0.8) = \qty{0.149}{L/min}$，\qty{8.9}{L}
每小时。
\textbf{15.} $0.83/(7\times 0.43) = \qty{0.28}{L/min}$：
\omterm{def:b1:gas-exchange:gill}{逆流}使需泵送的水量减半，做功亦然。
\textbf{16.} 每小时水量 \qty{8.9}{kg}；人每小时移动 $6\times
60\times 1.2 = \qty{430}{g}$ 的空气，质量仅为二十分之一，
而氧气量却多出三百倍（$15\,000$
对 \qty{50}{mL} 每小时）。
\textbf{17.} $100\times(0.95 - 0.15) = \qty{80}{mL/L}$；
心输出量 $0.83/0.080 = \qty{10.4}{mL/min}$。
\textbf{18.} 水/血 $149/10.4 = 14$；空气/血 $4.2/5 = 0.84$。
鱼须让十四倍于血液体积的水流过，
因为水的含氧量很低；人则约为一。
\textbf{19.} 摄氧量 $\times 1.5$，而水中含氧量仅为原来的
$5.5/7 = 0.79$ 倍：通气量 $\times 1.9$；呼吸的代价大致与流量成正比（
或更多，因其增长快于线性），在增大的摄氧量中所占比例由 \qty{10}{\%}
升至至少 \qty{13}{\%}，
鱼在温水中呼吸更费力。
\textbf{20.} 水入口端：$21 - 18 = \qty{3}{kPa}$；水出口
端：$4 - 3 = \qty{1}{kPa}$；两端梯度均为正。
\textbf{21.} 入口端 $21 - 3 = \qty{18}{kPa}$；出口端 $12 - 12 =
0$：两液体已达平衡，不再有梯度，尽管
水中仍保留其 \qty{57}{\%} 的氧。
\textbf{22.} 并流时，血液与单一、混合充分的\omterm{def:b1:gas-exchange:lung}{肺泡}气达到平衡，
并以该气体的压强离开；血液压不可超过
\qty{13.3}{kPa}，因为下游没有可相遇的更新鲜气体，
\omterm{def:b1:gas-exchange:gill}{逆流}中则存在这种情况。
\textbf{23.} 鳟鱼 $2000/(2\times 10^{-4}) = \num{1e7}$（cm$^2$ 每 cm）；
人 $10^6/(5\times 10^{-5}) = \num{2e10}$：比值 \num{2000}。
摄氧量之比为 $250/0.83 = 300$。鳟鱼每单位
$A/x$ 摄取的氧气多七倍：其跨鳃小片的平均压差必须更大
——这由\omterm{def:b1:gas-exchange:gill}{逆流}提供，使梯度沿整个
表面始终存在，而潮式肺的梯度在接近平衡
末端时已很小。
\textbf{24.} $\mathrm{CO_2}$ 在水中的溶解度
比氧气高二十五倍：水带走它的速度与产生速度相同，
故鱼血中的 $\mathrm{CO_2}$ 保持很低。
\textbf{25.} 鳃约占 \qty{80}{\%}，肺约占 \qty{20}{\%}
吸入介质中氧气的份额；鳃的单向水流及其
水与血液的\omterm{def:b1:gas-exchange:gill}{逆流}排列造成了这一差异。
\end{solution}
